\chapter{Geometric slices and their boundary identity}\label{sec:geometric-current-slices}
We separate the geometric level set, the distributional boundary, and the argument identifying them outside one null set. This supplies the boundary-identification argument used in Theorem \ref{thm:integral-current-slicing}. Rectifiable coarea, the countable $C^1$ decomposition, and the almost-everywhere fiber tangent assertion are now developed locally, with locally proved Euclidean area/coarea and Whitney extension retained as inputs. These are the geometric ingredients in \cite{Sim18}, Chapter 6, Lemmas 4.1 and 4.5, pp.~180--183.

\begin{lemma}[Orientation of a codimension-one fiber]\label{lem:slice-unit-orientation}
Let $P$ be an oriented $m$-plane, $m\geq1$, with unit simple orientation $\tau$, and let $v\in P$ be a unit vector. Define $\tau\mathbin{\lrcorner}v$ by
\[
 \langle\tau\mathbin{\lrcorner}v,\omega\rangle
 =\langle\tau,v^\flat\wedge\omega\rangle.
\]
It is the unit simple orientation of $P\cap v^\perp$ characterized by $v\wedge(\tau\mathbin{\lrcorner}v)=\tau$. When $m=1$, it is the signed unit zero-vector. This is the linear algebra in \cite{Sim18}, equation (4.3).
\end{lemma}
\begin{proof}
For $m\geq2$, choose an oriented orthonormal basis $(v,e_2,\ldots,e_m)$ of $P$. Then $\tau=v\wedge e_2\wedge\cdots\wedge e_m$, and evaluation on basis forms shows that contraction is $e_2\wedge\cdots\wedge e_m$. For $m=1$, write $\tau=\sigma v$ with $\sigma\in\{1,-1\}$. The contraction is the scalar $\sigma\in\Lambda^0P$, which gives the same identity and unit norm. These formulas give simplicity, norm one, the tangent plane, and the asserted orientation identity.
\end{proof}

\begin{definition}[The geometric level current]\label{def:geometric-level-current}
Let $T=\llbracket M,\theta,\tau\rrbracket$ be a locally integer rectifiable $m$-current of locally finite mass in an open set $U$, $m\geq1$, and let $f:U\to\mathbb R$ be locally Lipschitz. Write
\[
 M_t=\{x\in M:f(x)=t,\ |\nabla^Mf(x)|>0\},\qquad
 S_t=\llbracket M_t,\theta,\tau\mathbin{\lrcorner}
             (\nabla^Mf/|\nabla^Mf|)\rrbracket.
\]
This is defined at the almost-everywhere levels where the fiber is rectifiable, its tangent plane is $T_xM\cap\ker d^Mf_x$, and its weighted mass is locally finite. Values on the exceptional levels are not specified by this definition. See \cite{Sim18}, Lemma 4.1 and Definition 4.4.
\end{definition}

\begin{proposition}[Mass and integration of geometric slices]\label{prop:geometric-slice-integration}
The geometric level currents exist for almost every $t$. For every Borel $E\subset U$,
\[
 \int_{\mathbb R}\|S_t\|(E)\,dt
 =\int_E|\nabla^Mf|\,d\|T\|.
\]
For a smooth compactly supported $(m-1)$-form $\omega$ and $\psi\in C_c^\infty(\mathbb R)$,
\[
 \int_{\mathbb R}\psi(t)S_t(\omega)\,dt
 =\int_M\psi(f)\langle\tau,d^Mf\wedge\omega\rangle\,d\|T\|.
\]
These are the mass and signed coarea identities underlying \cite{Sim18}, Lemma 4.5(1).
\end{proposition}
\begin{proof}
For $m\geq2$, use weighted rectifiable coarea, Corollary \ref{cor:weighted-rectifiable-coarea}. For $m=1$, use Proposition \ref{prop:curve-scalar-coarea}. For $m\geq2$, the almost-everywhere tangent-plane assertion is proved in Theorem \ref{thm:rectifiable-coarea-chart-proof}, corresponding to the fiber input in \cite{Sim18}, Lemma 4.1. In dimension one a locally finite fiber consists of isolated points and has the zero-dimensional tangent. Lemma \ref{lem:slice-unit-orientation} makes the prescribed orientation unit and simple. The multiplicity is the positive integer $\theta$, so the mass measure is exactly $\theta\mathcal H^{m-1}\llcorner M_t$. Coarea with weight $\theta\chi_E$ proves the mass identity, including infinite values. With $E$ in a compact exhaustion, local Lipschitz bounds and local mass finiteness show that the slice has locally finite mass for almost every $t$, simultaneously on all compact subsets.

For the signed identity, insert the weight
\[
 \theta\psi(f)\left\langle\tau\mathbin{\lrcorner}
       \frac{\nabla^Mf}{|\nabla^Mf|},\omega\right\rangle
\]
on the noncritical part and set it to zero elsewhere. Its absolute value after multiplication by $|\nabla^Mf|$ is integrable on the compact support of $\omega$. Signed coarea and the contraction identity therefore give the formula. Critical points contribute zero almost everywhere in the level, so this convention causes no missing term.
\end{proof}

\begin{lemma}[An integrated identity for the distributional slice]\label{lem:integrated-distributional-slice}
Assume also that $T$ is locally normal, and put $B_t=\partial(T\llcorner\{f<t\})-(\partial T)\llcorner\{f<t\}$. For every test $(m-1)$-form $\omega$ and $\psi\in C_c^\infty(\mathbb R)$,
\[
 \int_{\mathbb R}\psi(t)B_t(\omega)\,dt
 =\int_{\mathbb R}\psi(t)S_t(\omega)\,dt.
\]
The left-hand side exists without assuming finite mass of $B_t$. This is an integrated version of the approximation argument in \cite{Sim18}, Lemma 4.5(2).
\end{lemma}
\begin{proof}
Set $Q(s)=\int_s^\infty\psi(t)\,dt$. Fubini, applied to the locally finite measures representing $T$ and $\partial T$ on the test support, gives
\[
 \int\psi(t)B_t(\omega)\,dt
 =(T\llcorner Q(f))(d\omega)
       -((\partial T)\llcorner Q(f))(\omega).
\]
The composition $Q(f)$ is locally Lipschitz. Proposition \ref{prop:lipschitz-rectifiable-product-rule} identifies this difference with $-(T\llcorner d^M(Q(f)))(\omega)$. Since $Q'=-\psi$, Lemma \ref{lem:tangential-scalar-chain} changes the negative sign to a positive one and gives
\[
 \int_M\psi(f)\langle\tau,d^Mf\wedge\omega\rangle\,d\|T\|.
\]
Proposition \ref{prop:geometric-slice-integration} identifies this with the required integral of $S_t$. The two scalar functions of $t$ are measurable and locally integrable: the first is bounded by the local masses times the $C^1$ norm of $\omega$, and the second is locally integrable by coarea.
\end{proof}

\begin{proposition}[A single exceptional null set for all test forms]\label{prop:slice-common-exceptional-set}
Under the preceding assumptions there is a Lebesgue-null set $N\subset\mathbb R$ such that
\[
 B_t=S_t\quad\text{as currents for every }t\notin N.
\]
Thus equality is not merely an almost-everywhere assertion with a different exceptional set for each test form. This expands the countable-density step in \cite{Sim18}, p.~183.
\end{proposition}
\begin{proof}
For a fixed test form $\omega$, Lemma \ref{lem:integrated-distributional-slice} says that the locally integrable function $B_t(\omega)-S_t(\omega)$ has zero pairing with every smooth compactly supported scalar function. Lemma \ref{lem:zero-distribution-function} makes it zero almost everywhere.

Choose countable $C^1$-dense test families on a relatively compact exhaustion of $U$ using Lemma \ref{lem:countable-current-c1-tests}. The approximants have a common slightly larger compact support. Remove the null exceptional set for each member of all these countable families, together with the set where $S_t$ fails to have locally finite mass.

For each remaining $t$, a test supported in a compact $K$ satisfies
\[
 |B_t(\omega)|\leq\|T\|(K)\|d\omega\|_{*,\infty}
                  +\|\partial T\|(K)\|\omega\|_{*,\infty},
 \qquad |S_t(\omega)|\leq\|S_t\|(K)\|\omega\|_{*,\infty}.
\]
Use the slightly larger compact set for the same estimates on approximating forms. $C^1$ convergence therefore extends equality from the chosen countable family to every test. The removed set is countable union of null sets, proving the claim. The $C^1$ topology is necessary here because the distributional expression for $B_t$ contains $d\omega$ before its finite-mass representation is established.
\end{proof}

\begin{proposition}[Integral restrictions and integral slice boundaries]\label{prop:integral-geometric-cut}
If $T$ is locally integral, then for almost every $t$ both $T\llcorner\{f<t\}$ and $S_t$ are locally integral. For $m\geq2$,
\[
 \partial S_t=-S_t^{\partial T},
\]
where $S_t^{\partial T}$ denotes the geometric slice constructed from $\partial T$. For $m=1$ the boundary of the zero-current $S_t$ is zero. See \cite{Sim18}, Lemma 4.5(2)--(3).
\end{proposition}
\begin{proof}
Apply Proposition \ref{prop:slice-common-exceptional-set} to $T$. If $m\geq2$, also apply it to $\partial T$, which is locally integer rectifiable and locally normal because $\partial^2T=0$. Intersect the two full-measure sets of good levels. Taking a distributional boundary in the identity for $S_t$ gives
\[
 \partial S_t=-\partial((\partial T)\llcorner\{f<t\})
             =-S_t^{\partial T}.
\]
The latter is integer rectifiable with locally finite mass. Thus $S_t$ and its boundary are integer rectifiable and locally finite, which is precisely local integrality. In dimension one its boundary vanishes by convention, giving the same conclusion directly. The restriction $T\llcorner\{f<t\}$ is integer rectifiable, and
\[
 \partial(T\llcorner\{f<t\})=(\partial T)\llcorner\{f<t\}+S_t
\]
has the same properties and locally finite mass. Hence the restriction is locally integral. No compactness theorem is needed for these deductions.
\end{proof}
